%% ============================================================================
%% Section I: Introduction  +  Section II: Related Work
%% ============================================================================

\section{Introduction}
\label{sec:intro}

\IEEEPARstart{L}{arge} language models (LLMs) excel at open-domain question
answering, but deploying them for a specific knowledge domain exposes three
persistent obstacles. First, commercial API pricing places a recurring cost
on every query, which is prohibitive for resource-constrained organizations
such as churches, seminaries, and non-profit publishers. Second, sensitive
corpora and user questions often cannot leave the premises, ruling out
cloud-only inference. Third, and most specific to the religious domain,
factual grounding is non-negotiable: recent measurements report that even
state-of-the-art LLMs score lowest on faith-related dimensions of a
human-flourishing alignment benchmark~\cite{hilliard2025flourishing},
reinforcing a legitimate concern in the Christian community about whether
LLM answers are actually grounded in Scripture rather than in the model's
parametric memory.

Retrieval-augmented generation (RAG)~\cite{lewis2020rag} mitigates factual
errors by conditioning generation on retrieved evidence, and graph-based
RAG (GraphRAG) variants~\cite{edge2024graphrag,peng2026graphragsurvey}
additionally exploit entity and relation structure for multi-hop and
aggregation queries. However, most published GraphRAG systems are evaluated
as end-to-end black boxes: when a knowledge graph (KG) improvement does or
does not move the final metric, the community rarely learns \emph{which
layer} --- extraction, storage, retrieval, ranking, generation, or the
benchmark itself --- gated the outcome. This paper documents a deployed
system in which exactly that question was forced into the open, three
times, by concrete engineering incidents.

We design, build, deploy, and iteratively optimize a GraphRAG system for
question answering over the Traditional Chinese Bible (Chinese Union
Version, new punctuation edition; CUV). The system converts 66 books from
PDF into a pericope-centric hierarchical corpus, constructs a knowledge
graph with grounded entity and relation extraction, indexes three
complementary databases (PostgreSQL, Qdrant, Neo4j), and serves queries
through a signal-driven six-route retrieval architecture terminated by a
cross-encoder reranker and a linear rank-fusion layer. Along the way we
performed a quantified audit of KG construction quality, executed a
completeness overhaul, observed a rigorously verified \emph{negative
result} --- the overhaul produced zero retrieval improvement on our
benchmark --- diagnosed the mechanism question by question, and repaired
the actual bottleneck (the ranking layer) in one day of engineering,
surpassing all previous results.

This paper makes six contributions:

\begin{enumerate}
\item \textbf{A complete, reproducible pipeline from raw PDF to GraphRAG
for Traditional Chinese biblical text} (Sections~\ref{sec:preprocessing}
and~\ref{sec:kg}): layout-geometry PDF parsing without OCR or LLMs,
pericope-centric hierarchical chunking with the embedding model's own
tokenizer, grounded dual-track entity extraction (six types, 9{,}120
entities, 173{,}896 mentions), closed-schema grounded relation extraction
(37 relation types with evidence-span verification), and a three-source
cross-reference network of 250{,}418 edges integrating the public-domain
Treasury of Scripture Knowledge (TSK)~\cite{tsk}.

\item \textbf{A signal-adaptive six-route retrieval architecture over a
triple-database backend, terminated by a rank-fusion layer}
(Section~\ref{sec:retrieval}): six boolean query signals select among six
routes plus a fallback; each route runs a weighted combination of eleven
retrieval strategies; candidates are deduplicated, reranked by a
cross-encoder, fused by
$s_{\mathrm{fused}}=(1-\alpha)\,s_{\mathrm{rr}}+\alpha\,w$, and finally
adjusted by a small set of ``user-literal'' pins.

\item \textbf{A quantified causal chain from chunking granularity to KG
quality} (Section~\ref{sec:audit}): we show that the damage to KG quality
was caused not by chunk-size parameters but by the \emph{silent
inheritance} of an embedding-oriented granularity structure by the
extraction, import, and retrieval layers (mechanisms M1--M6), producing a
55.9\% silent mention loss and 1{,}474 retrieval-blind entities. The
finding is aligned with, and locally re-validates, published evidence on
extraction granularity~\cite{edge2024graphrag}, document-level relation
extraction~\cite{yao2019docred}, and coreference-resolved KG
construction~\cite{meher2025corekg,meher2025insidecorekg,meher2025linkkg}.

\item \textbf{A verified negative result with a mechanism-level diagnosis}
(Section~\ref{sec:p0}): a KG completeness overhaul that raised entity
anchor coverage from 83.8\% to 98.2\% and multiplied cross-references by
273 produced \emph{no} improvement --- and locally a regression --- on a
100-question head-entity benchmark. Question-by-question forensics
attribute this to (a) topical cross-reference votes acting as noise on
event-narrative questions, (b) a final ranking monopolized by the
cross-encoder's surface-form score, and (c) benchmark head saturation. This
result complements benchmark-level observations that KG defects propagate
non-linearly to answers~\cite{zhou2025brink} and that GraphRAG gains are
frequently over-read due to evaluation bias~\cite{zeng2025unbiased}.

\item \textbf{A rank-fusion remediation with an interpretable ablation}
(Section~\ref{sec:fusion}): TSK de-noising by vote-weighted edge
declassing, a curated head-event backfill that re-enabled entity-vector
retrieval, and the linear fusion layer together lift hit rate to 0.97
(event-narrative questions to 1.00), surpassing the pre-overhaul state and
widening the margin over the semantic-only baseline to $+0.16$. The
$\alpha$ ablation exposes a complementary structure --- discrete data
repairs saturate hit rate, continuous fusion recovers tail anchors --- and
a mechanism-level finding that discrete ``pin'' patches flip from
net-positive to net-negative once continuous fusion exists.

\item \textbf{A methodological conclusion for GraphRAG practice}
(Section~\ref{sec:discussion}): KG construction quality, the retrieval
ranking mechanism, and benchmark design jointly determine the observable
benefit of any single-layer investment. In our deployment, one day of
ranking-layer repair outperformed what a week-scale re-extraction was
projected to deliver on the same benchmark; three independent incident
rounds show graph assets whose value was gated by retrieval-side
mechanisms. We argue that evaluation design is a \emph{prerequisite} for
--- not an afterthought of --- further construction-side optimization.
\end{enumerate}

The remainder of the paper follows the pipeline order: corpus and
preprocessing (Section~\ref{sec:preprocessing}), knowledge-graph
construction (Section~\ref{sec:kg}), online retrieval and generation
(Section~\ref{sec:retrieval}), the evaluation framework
(Section~\ref{sec:eval}), the quality audit and the three experiment
rounds (Sections~\ref{sec:round0}--\ref{sec:fusion}), discussion
(Section~\ref{sec:discussion}), limitations and future work
(Sections~\ref{sec:limitations}--\ref{sec:future}), and conclusion
(Section~\ref{sec:conclusion}). Appendix~\ref{app:citations} maps each
pipeline stage to the literature it draws on.

\section{Related Work}
\label{sec:related}

\subsection{Retrieval-Augmented Generation and Hybrid Retrieval}

RAG~\cite{lewis2020rag} couples a parametric generator with non-parametric
retrieval; surveys~\cite{gao2023ragsurvey} catalogue its now-extensive
design space. Dense passage retrieval~\cite{karpukhin2020dpr} and lexical
scoring~\cite{robertson2009bm25} remain the two complementary retrieval
substrates, commonly combined by reciprocal rank fusion
(RRF)~\cite{cormack2009rrf}. Our system uses BGE-M3~\cite{chen2024bgem3}
for dense embeddings (and its companion cross-encoder,
BGE-reranker-v2-m3, for reranking), BM25 with Chinese-specific CKIP
tokenization~\cite{ckip} for sparse vectors, and RRF for dense--sparse
hybrid search. Distinct from standard hybrid search, the final ranking in
our system linearly fuses the cross-encoder score with a
\emph{strategy-prior} weight that encodes which retrieval strategy
proposed the candidate --- a channel by which graph evidence survives into
the final top-$k$ (Section~\ref{sec:fusionlayer}). Calibrated fusion of
heterogeneous graph and vector scores has recently been studied in
multi-hop QA~\cite{bacellar2026calibrated}; our fusion is deliberately
simpler (a single linear interpolation with one global $\alpha$) and is
evaluated by ablation.

\subsection{Graph-Based RAG}

GraphRAG~\cite{edge2024graphrag} popularized LLM-built entity graphs with
hierarchical community summaries for query-focused summarization;
subsequent systems explore lighter indexing and different retrieval
operators, \eg LightRAG's dual-level keyword retrieval~\cite{guo2024lightrag},
personalized-PageRank retrieval over passage-linked
graphs~\cite{gutierrez2025hipporag2}, attributed-community hierarchical
retrieval~\cite{wang2025archrag}, and recursive abstractive summary trees
(RAPTOR)~\cite{sarthi2024raptor}. Event-centric variants build explicit
event nodes and temporal edges~\cite{zhang2025entityevent,sun2025dygrag}.
Benchmark studies ask \emph{when} graphs actually help
RAG~\cite{xiang2025whengraphs}. Our architecture differs from these
retrieval-operator designs in that graph traversal is one of ten
strategies selected \emph{per query} by boolean signals (route-level
adaptivity), and in that the paper's empirical focus is not a new
operator but the interaction between graph asset quality, ranking, and
benchmark --- an interaction that, to our knowledge, has not been
documented at incident-level granularity for a deployed GraphRAG system.

\subsection{LLM-Based Knowledge Graph Construction}

LLM-driven KG construction is surveyed in~\cite{bian2025kgsurvey}, which
contrasts schema-based and schema-free extraction paradigms; our relation
extractor is schema-based with a closed 37-type ontology. On construction
quality, GraphRAG's own ablation shows chunk granularity directly controls
extraction density (600- vs.\ 2{,}400-token chunks yield nearly twice the
entity references)~\cite{edge2024graphrag}; DocRED quantifies that at
least 40.7\% of relational facts require cross-sentence evidence and
17.6\% require coreference reasoning~\cite{yao2019docred}; and
``lost-in-the-middle'' effects bound how much context can be packed into a
single extraction call~\cite{liu2024lost}. Coreference-resolved
construction reduces node duplication by 28--45\% in the CORE-KG /
LINK-KG line of work~\cite{meher2025corekg,meher2025insidecorekg,%
meher2025linkkg}, with the ablation in~\cite{meher2025insidecorekg}
cleanly separating coreference (deduplication) from structured prompting
(noise suppression). Post-extraction correction is an alternative to
re-extraction: ontology-grounded repair~\cite{loconte2026betterlater},
canonicalization after open extraction~\cite{zhang2024edc}, KG denoising
that \emph{shrinks} graphs yet improves downstream RAG~\cite{zheng2025lessismore},
retrieval-augmented construction that re-retrieves cross-chunk context for
each pre-entity~\cite{zhang2025rakg}, restoring the source context of
isolated triples~\cite{huang2025tcrqf}, and GNN-guided cross-chunk graph
augmentation~\cite{zhang2026crossaug}. Atomic-fact decomposition with
dual-time modeling~\cite{lairgi2025atom} shows finer construction units
increase fact coverage. Our audit (Section~\ref{sec:audit}) maps each of
our locally measured defects onto this literature, and our remediation
roadmap adopts its recipes.

\subsection{Evaluating RAG and GraphRAG}

We evaluate generation with RAGAS~\cite{es2023ragas} plus a custom
LLM-judged answer-coverage metric, and retrieval with deterministic
ground-truth metrics. Two recent results shape our methodology. BRINK
benchmarks how KG-based RAG breaks under incomplete knowledge and finds
that answer quality degrades non-linearly with KG defects, with systems
partly relying on parametric memory rather than structural
reasoning~\cite{zhou2025brink}. An unbiased evaluation framework for
GraphRAG shows that after removing position and verbosity bias from
LLM-judged comparisons, reported win rates can drop by tens of points
(\eg LightRAG $66.70\%\!\rightarrow\!39.06\%$)~\cite{zeng2025unbiased}. Our
own evaluation encountered a concrete instance of judge bias --- the RAGAS
answer-relevancy ``noncommittal'' rule zeroes honest refusals, deflating
the score of a cautious commercial model (Section~\ref{sec:evalmetrics}).
Chunking-strategy sensitivity for \emph{retrieval} (not construction) has
been studied systematically in~\cite{shaukat2026chunking}; we cite it only
for its retrieval-side scope, since no peer-reviewed quantification exists
for overlap effects on \emph{extraction}, and we rely on
DocRED~\cite{yao2019docred} as indirect evidence there.
